% docs/manual/operations/chapters/07-monitoring.tex
% 第 7 章：监控与日志

\chapter{监控与日志}

长跑任务必须能监控。本章详解 \openbench{} 的日志系统、进度事件、报告生成，以及如何接入外部监控（Prometheus / 日志聚合）。

\section{logging\_system 总览}

\file{util/logging\_system.py}（581 行）配置 Python 标准 logging：

\begin{itemize}
  \item Console handler：默认输出到 stderr
  \item File handler：写到 \file{output/<run>/run.log}
  \item Format：时间 + 模块名 + 级别 + 消息
\end{itemize}

每个模块用 \texttt{logging.getLogger(\_\_name\_\_)} 拿 logger，命名形如 \texttt{openbench.config.loader}、\texttt{openbench.runner.local}。

\section{日志级别配置}

\subsection{在 yaml 中}

\begin{minted}{yaml}
project:
  debug_mode: true   # DEBUG 级别（详细）；默认 INFO
\end{minted}

只有总开关。子模块级控制要程序化：

\begin{minted}{python}
import logging
logging.getLogger("openbench.runner.local").setLevel(logging.DEBUG)
logging.getLogger("openbench.data.pipeline").setLevel(logging.WARNING)
\end{minted}

\subsection{级别约定}

详见开发卷附录 E。运维角度速查：

\begin{itemize}
  \item DEBUG：内部细节；只看
  \item INFO：阶段切换、cache hit、output written；正常运行的"心跳"
  \item WARNING：deprecation、低置信元信息、自动 fallback；\textbf{建议告警}
  \item ERROR：单 task 失败但程序继续；\textbf{必须告警}
  \item CRITICAL：致命，程序退出；\textbf{紧急告警}
\end{itemize}

\section{日志去向}

\subsection{stderr}

默认 \openbench{} 把日志送 stderr。在终端能看到。命令行处置：

\begin{minted}{bash}
openbench run config.yaml 2> run.err           # 只重定向 stderr
openbench run config.yaml 2>&1 | tee all.log   # 合并到文件 + 屏幕
\end{minted}

\subsection{output/<run>/run.log}

每次 \cli{run} 自动写完整日志到这里。后处理友好：

\begin{minted}{bash}
# 看末尾
tail -100 output/<run>/run.log

# 实时跟
tail -f output/<run>/run.log

# 找错误
grep -E "ERROR|CRITICAL" output/<run>/run.log
\end{minted}

\subsection{大日志的处理}

debug\_mode + 长任务可能产出 GB 级日志。\openbench{} 没自动 rotation，需手动：

\begin{minted}{bash}
# 实时压缩归档
gzip output/<run>/run.log
# 或定期 logrotate（管理员配 /etc/logrotate.d/openbench）
\end{minted}

\section{progress.py：进度事件}

\file{util/progress.py} 把 task 状态封装为可解析的进度事件，供 GUI 与 CLI 共用：

\subsection{事件格式}

INFO 级别消息中可识别的进度模式：

\begin{minted}{text}
[INFO] Phase: preprocess
[INFO] Queued Evapotranspiration: sim=CoLM2024 ref=GLEAM_v4.2a
[INFO] Cache hit: skipping <task_key>
[INFO] Output written: output/<run>/metrics/<var>/<sim>_metrics.csv
[INFO] Phase: comparison
\end{minted}

\subsection{用作进度条}

GUI 用 \modname{gui.progress\_parser} 解析这些事件更新进度。CLI 用户可自己 grep：

\begin{minted}{bash}
# 简单计数
grep -c "Output written" output/<run>/run.log
\end{minted}

\section{report.py：HTML/PDF 总结}

\file{util/report.py} 在评估末尾生成报告。运维相关：

\subsection{报告生成的开销}

\begin{itemize}
  \item HTML：~5-15 秒（含模板渲染 + 图嵌入）
  \item PDF：额外 30-60 秒（要 \cli{[report]} extra 装 \modname{xhtml2pdf}）
\end{itemize}

赶时间可关：

\begin{minted}{yaml}
project:
  generate_report: false
\end{minted}

\subsection{自定义模板}

\modname{report} 用 \modname{jinja2} 模板。模板路径在包内（不容易改）。要自定义：

\begin{itemize}
  \item Fork 仓库改 \file{util/report.py} 内嵌的模板
  \item 或评估完后写自己的脚本聚合 CSV/figure 出报告
\end{itemize}

\section{集成外部监控}

\openbench{} 没内置 Prometheus / DataDog / 等 exporter，但日志结构允许接入。

\subsection{方案 A：日志聚合 + 告警}

部署 Loki / ELK / Splunk 等，把 \openbench{} 日志接入。在聚合系统设：

\begin{itemize}
  \item ERROR 出现 → 告警
  \item WARNING 频次超阈值 → 告警
  \item INFO "Phase" 长时间不切换 → 卡住告警
\end{itemize}

\subsection{方案 B：脚本采样}

简单做法：cron 每 5 分钟读 \file{run.log}，提取统计：

\begin{minted}{bash}
#!/bin/bash
LOG=output/<run>/run.log
TOTAL_TASKS=$(grep -c "^\[INFO\] Queued" $LOG)
DONE_TASKS=$(grep -c "Output written" $LOG)
ERRORS=$(grep -c "^\[ERROR\]" $LOG)
echo "$(date) progress=$DONE_TASKS/$TOTAL_TASKS errors=$ERRORS"
\end{minted}

\subsection{方案 C：自定义 logging handler}

在脚本中插入 Prometheus pushgateway handler：

\begin{minted}{python}
import logging
from prometheus_client import push_to_gateway, Counter, REGISTRY

errors = Counter('openbench_errors', 'OpenBench errors')

class PrometheusHandler(logging.Handler):
    def emit(self, record):
        if record.levelno >= logging.ERROR:
            errors.inc()
            push_to_gateway('localhost:9091', job='openbench', registry=REGISTRY)

logging.getLogger("openbench").addHandler(PrometheusHandler())
\end{minted}

之后照常调 \texttt{run\_evaluation()}。

\section{长任务监控建议}

任务跑 > 1 小时建议：

\begin{enumerate}
  \item SLURM job 用 \texttt{--mail-type=END,FAIL --mail-user=you@example.com}
  \item 自定义脚本完成时发 Slack/email 提醒
  \item 在 dashboard 显示进度（grep 日志的简单方案就够）
  \item 失败时立即停下不要等到全完才发现
\end{enumerate}

\section{日志归档与清理}

\openbench{} 不自动归档/清理日志。建议管理员做：

\begin{minted}{bash}
# /etc/cron.daily/openbench-log-rotate
find /shared/openbench -name "run.log" -mtime +30 -exec gzip {} \;
find /shared/openbench -name "run.log.gz" -mtime +180 -delete
\end{minted}

\section{下一步}

下一章详解故障恢复：中断、断电、磁盘满、死锁等的诊断与恢复流程。
